\choiceheading

\begin{question}[index=1]
设 $f(x)$ 为微分方程 $y'+\left(x+\dfrac12\right)y=\displaystyle\int_0^1 |x-t|\dd t$
满足 $y(0)=1$ 的解，则有\paren
\begin{choices}[columns=1]
\item 在点 $x=0$ 处 $f(x)$ 取极大值
\item 在点 $x=0$ 处 $f(x)$ 取极小值
\item 点 $(0,f(0))$ 为曲线 $y=f(x)$ 的拐点
\item 点 $x=0$ 不是 $f(x)$ 的极值点，点 $(0,f(0))$ 也不是曲线 $y=f(x)$ 的拐点
\end{choices}
\end{question}

\begin{question}[index=2]
设 $f(x,y)$ 与 $\varphi(x,y)$ 均为可微函数，且 $\varphi'_y(x,y)\ne 0$。
已知 $(x_0,y_0)$ 是 $f(x,y)$ 在约束条件 $\varphi(x,y)=0$ 下的一个极值点，下列选项正确的是\paren
\begin{choices}[columns=2]
\item 若 $f'_x(x_0,y_0)=0$，则 $f'_y(x_0,y_0)=0$
\item 若 $f'_x(x_0,y_0)=0$，则 $f'_y(x_0,y_0)\ne 0$
\item 若 $f'_x(x_0,y_0)\ne 0$，则 $f'_y(x_0,y_0)=0$
\item 若 $f'_x(x_0,y_0)\ne 0$，则 $f'_y(x_0,y_0)\ne 0$
\end{choices}
\end{question}

\begin{question}[index=3]
设 $J_i=\displaystyle\iint_{D_i}\sqrt[3]{x-y}\dd x\dd y$（$i=1,2,3$），其中
\[
\begin{gathered}
D_1=\{(x,y)\mid 0\le x\le 1,\ 0\le y\le 1\},\\
D_2=\{(x,y)\mid 0\le x\le 1,\ 0\le y\le\sqrt{x}\},\qquad
D_3=\{(x,y)\mid 0\le x\le 1,\ x^2\le y\le 1\},
\end{gathered}
\]
则\paren
\begin{choices}[columns=4]
\item $J_1<J_2<J_3$
\item $J_3<J_1<J_2$
\item $J_2<J_3<J_1$
\item $J_2<J_1<J_3$
\end{choices}
\end{question}

\begin{question}[index=4]
设函数 $f(x)$ 为 $[0,1]$ 上的可导函数，且 $|f'(x)|$ 在 $[0,1]$ 上的最大值为 $M$。
若方程 $f(x)-x+1=0$，$f(x)+x-1=0$ 在 $(0,1)$ 内均有解，则对于满足条件的函数 $f(x)$，均有\paren
\begin{choices}[columns=2]
\item $|f(0)|+|f(1)|\le M$
\item $|f(0)|+|f(1)|\ge M$
\item $|f(0)|+|f(1)|=M$
\item $|f(0)|+|f(1)|\ne M$
\end{choices}
\end{question}

\begin{question}[index=5]
设 $\symbfit{A}$ 与 $\symbfit{B}$ 都是 $n$ 阶方阵，$\symbfit{A}\symbfit{x}=\symbf{0}$ 的解都是
$\symbfit{B}\symbfit{x}=\symbf{0}$ 的解，$\symbfit{\xi}_1,\symbfit{\xi}_2,\symbfit{\xi}_3$ 为
$\symbfit{A}\symbfit{x}=\symbf{0}$ 的基础解系，则在下列线性方程中以
$\symbfit{\xi}_1+\symbfit{\xi}_2,\symbfit{\xi}_2+\symbfit{\xi}_3,\symbfit{\xi}_1+\symbfit{\xi}_3$ 为基础解系的是\paren
\begin{choices}[columns=4]
\item $(\symbfit{A}+\symbfit{B})\symbfit{x}=\symbf{0}$
\item $(\symbfit{A}-\symbfit{B})\symbfit{x}=\symbf{0}$
\item $\symbfit{A}\symbfit{B}\symbfit{x}=\symbf{0}$
\item $\begin{pmatrix}\symbfit{A}\\\symbfit{B}\end{pmatrix}\symbfit{x}=\symbf{0}$
\end{choices}
\end{question}

\begin{question}[index=6]
设向量
\[
\symbfit{\alpha}_1=\begin{bmatrix}a\\1\\-1\\1\end{bmatrix},\qquad
\symbfit{\alpha}_2=\begin{bmatrix}1\\1\\b\\a\end{bmatrix},\qquad
\symbfit{\alpha}_3=\begin{bmatrix}1\\a\\-1\\1\end{bmatrix}.
\]
若 $\symbfit{\alpha}_1,\symbfit{\alpha}_2,\symbfit{\alpha}_3$ 线性相关，且其中任意两个向量均线性无关，则\paren
\begin{choices}[columns=4]
\item $a=1,\ b\ne -1$
\item $a=1,\ b=-1$
\item $a\ne -2,\ b=2$
\item $a=-2,\ b=2$
\end{choices}
\end{question}

\begin{question}[index=7]
设 $\symbfit{A}$ 为 $3$ 阶矩阵，$3$ 维列向量 $\symbfit{\alpha}_1,\symbfit{\alpha}_2,\symbfit{\alpha}_3$ 线性无关，
$\symbfit{A}(\symbfit{\alpha}_1,\symbfit{\alpha}_2,\symbfit{\alpha}_3)
=(\symbfit{\alpha}_1,\symbfit{\alpha}_2,-\symbfit{\alpha}_3)$，记
$\symbfit{P}=(\symbfit{\alpha}_1+\symbfit{\alpha}_3,\symbfit{\alpha}_2,-\symbfit{\alpha}_3)$，
则 $\symbfit{P}^{-1}\symbfit{A}\symbfit{P}=$\paren
\begin{choices}[columns=4]
\item $\begin{bmatrix}1&0&0\\0&1&0\\2&0&-1\end{bmatrix}$
\item $\begin{bmatrix}1&0&0\\0&1&0\\0&0&-1\end{bmatrix}$
\item $\begin{bmatrix}1&0&0\\0&1&0\\1&0&-1\end{bmatrix}$
\item $\begin{bmatrix}1&0&0\\0&1&0\\-1&0&1\end{bmatrix}$
\end{choices}
\end{question}

\begin{question}[index=8]
设随机变量 $(X,Y)$ 服从二维正态分布 $N\left(0,0;1,4;-\dfrac12\right)$，则下列随机变量中服从标准正态分布且与 $X$ 独立的是\paren
\begin{choices}[columns=4]
\item $\dfrac{\sqrt{5}}{5}(X+Y)$
\item $\dfrac{\sqrt{5}}{5}(X-Y)$
\item $\dfrac{\sqrt{3}}{3}(X+Y)$
\item $\dfrac{\sqrt{3}}{3}(X-Y)$
\end{choices}
\end{question}

\begin{question}[index=9]
已知随机变量 $X\dist N(0,1)$，$Y\dist B\left(1,\dfrac12\right)$，且 $X$ 和 $Y$ 相互独立，记 $Z_1=X+Y$，$Z_2=X\cdot Y$，则下列结论正确的是\paren
\begin{choices}[columns=1]
\item $Z_1,Z_2$ 均为非连续型随机变量
\item $Z_1$ 为非连续型随机变量，$Z_2$ 为连续型随机变量
\item $Z_1$ 为连续型随机变量，$Z_2$ 为非连续型随机变量
\item $Z_1,Z_2$ 均为连续型随机变量
\end{choices}
\end{question}

\begin{question}[index=10]
设总体 $X$ 的概率分布为
\[
\setlength{\arraycolsep}{13pt}
\renewcommand{\arraystretch}{1.45}
\begin{array}{c|ccc}
\hline
X & 1 & 2 & 3\\ \hline
P & \frac12 & \frac13 & \frac16\\ \hline
\end{array}
\]
以 $Y_i$ 表示来自总体 $X$ 的简单随机样本（样本容量为 $n$）中等于 $i$ 的个数（$i=1,2,3$），则 $Y_1$ 和 $Y_2$ 的相关系数是\paren
\begin{choices}[columns=4]
\item $\dfrac13$
\item $-\dfrac{\sqrt{3}}{3}$
\item $-\dfrac{\sqrt{2}}{2}$
\item $-\dfrac12$
\end{choices}
\end{question}

\fillheading

\begin{question}[index=11]
设 $L$ 是空间圆周 $\begin{cases}
x^2+y^2+z^2=a^2,\\
x+y+z=\dfrac32a,
\end{cases}$ 则
$\displaystyle\oint_L(2yz+2zx+2xy)\dd s=$\fillin。
\end{question}

\begin{question}[index=12]
已知函数 $f(x)$ 在 $(-\infty,+\infty)$ 上连续，且 $f(x)=(x+1)^2+2\displaystyle\int_0^x f(t)\dd t$，
则当 $n\ge 2$ 时，$f^{(n)}(0)=$\fillin。
\end{question}

\begin{question}[index=13]
方程 $x\dfrac{\mathrm{d}y}{\mathrm{d}x}+y=\dfrac{1}{1+(xy)^2}$ 的通解为\fillin。
\end{question}

\begin{question}[index=14]
极限
$\displaystyle\lim_{n\to\infty}\sum_{i=1}^{n}\left[\frac{1}{(n+i+1)^2}+\frac{1}{(n+i+2)^2}+\cdots+\frac{1}{(n+i+i)^2}\right]=$\fillin。
\end{question}

\Needspace{6\baselineskip}
\begin{question}[index=15]
设有 $n$ 元实二次型
\[
f(x_1,x_2,\ldots,x_n)=(x_1+a_1x_2)^2+(x_2+a_2x_3)^2+\cdots
+(x_{n-1}+a_{n-1}x_n)^2+(x_n+a_nx_1)^2,
\]
其中 $a_i$（$i=1,2,\ldots,n$）为实数。\par
该二次型为正定二次型时，$a_1,a_2,\ldots,a_n$ 应满足的条件为\fillin。
\end{question}

\begin{question}[index=16]
设总体 $X$ 的概率密度为 $f(x)=\dfrac{1}{2\theta}\eu^{-\frac{|x|}{\theta}}$，$-\infty<x<+\infty$，$\theta>0$。
$X_1,X_2,\ldots,X_n$ 是取自总体 $X$ 的样本，则未知参数 $\theta$ 的矩估计量为\fillin。
\end{question}

\solveheading

\begin{question}[index=17,points=10,show-points=true,points-prelabel={（本题满分},points-postlabel={分）},points-separate-par=false]
已知函数 $f(x)$ 满足 $f(x)=\dfrac{1}{(2-x)^2}-\displaystyle\int_0^1 f(x)\dd x$，\\
将 $f(x)$ 展开成 $x$ 的幂级数。
\end{question}

\begin{question}[index=18,points=12,show-points=true,points-prelabel={（本题满分},points-postlabel={分）},points-separate-par=false]
设函数 $f(x,y)$ 具有连续的偏导数且满足 $\dfrac{f'_x(x,y)}{f(x,y)}=\dfrac{1}{x-y}-2x$，\\
以及 $f(x,-x)=2x\eu^{-2x^2}$。
\begin{enumerate}[label=（\Roman*）,leftmargin=2.6em,labelsep=.4em,itemsep=3pt,topsep=3pt]
\item 求 $f(x,y)$；
\item 求 $f(x,y)$ 的极值。
\end{enumerate}
\end{question}

\begin{question}[index=19,points=12,show-points=true,points-prelabel={（本题满分},points-postlabel={分）},points-separate-par=false]
计算 $I=\displaystyle\iiint_{\Omega}\dfrac{z}{\sqrt{x^2+y^2}}\dd v$，其中 $\Omega$ 是由 $yOz$ 平面上的区域 $D$（$y^2+z^2\le 5$，$z\ge y-1$，$y\ge 0$，$z\ge 0$）绕 $z$ 轴旋转一周所得的空间区域。
\end{question}

\begin{question}[index=20,points=12,show-points=true,points-prelabel={（本题满分},points-postlabel={分）},points-separate-par=false]
设 $x_1=a\ge 0$，$y_1=b\ge 0$，$a\le b$，$x_{n+1}=\sqrt{x_ny_n}$，
$y_{n+1}=\dfrac{x_n+y_n}{2}$（$n=1,2,\ldots$），证明 $\displaystyle\lim_{n\to\infty}x_n=\lim_{n\to\infty}y_n$。
\end{question}

\begin{question}[index=21,points=12,show-points=true,points-prelabel={（本题满分},points-postlabel={分）},points-separate-par=false]
设二次型 $f(x_1,x_2)=x_1^2-4x_1x_2+4x_2^2$ 经正交变换
$\begin{pmatrix}x_1\\x_2\end{pmatrix}=\symbfit{Q}\begin{pmatrix}y_1\\y_2\end{pmatrix}$\\
化为二次型 $g(y_1,y_2)=ay_1^2+4y_1y_2+by_2^2$，其中 $a\ge b$。
\begin{enumerate}[label=（\Roman*）,leftmargin=2.6em,labelsep=.4em,itemsep=3pt,topsep=3pt]
\item 求 $a,b$ 的值；
\item 求正交矩阵 $\symbfit{Q}$。
\end{enumerate}
\end{question}

\begin{question}[index=22,points=12,show-points=true,points-prelabel={（本题满分},points-postlabel={分）},points-separate-par=false]
两台仪器独立测试：仪器甲单次达标概率 $p=\dfrac13$，$X$ 是甲第一次达标所需测试次数；
仪器乙单次达标概率 $r=\dfrac12$，$Y$ 是乙第一次达标所需测试次数，$X,Y$ 相互独立。
\begin{enumerate}[label=（\arabic*）,leftmargin=2.6em,labelsep=.4em,itemsep=3pt,topsep=3pt]
\item 已经对甲连续测试 $6$ 次全部不达标，求至少再测试 $5$ 次甲才第一次达标的概率；
\item 计算期望 $E(X\mid X>6)$；
\item 判断 $W=\min\{X,Y\}$ 是否具备无记忆性，并给出证明。
\end{enumerate}
\end{question}
